% =====================================================================
%  24.3  The centroid tree as a distance index -- one table per centroid.
%  Exports: sec:c18-index  def:c18-index  thm:c18-nn  lst:c18-index
%           prob:c18-293e  lst:c18-293e
%  Solved CF 293E "Close Vertices" (2700), live-verified via the problemset
%  API (tags data structures / divide and conquer / trees).  Statement:
%  count unordered pairs (u,v) whose UNIQUE path uses at most l edges AND has
%  total edge weight at most w.  Official samples read live and replayed:
%    S1: n=4, l=4, w=6, star centred at 1 (edges 1-2=3,1-3=4,1-4=3) -> 4.
%        Centroid 1 claims: (1,2),(1,3),(1,4),(2,4); (2,3),(3,4) break w. -> 4.
%    S2: n=6, l=2, w=17, tree 1-2(3),2-3(5),2-4(13),1-5(6),5-6(9) -> 9.
%        Top centroid 1 claims 6 pairs; recursion on {2,3,4} claims 2 and on
%        {5,6} claims 1;  6+2+1 = 9.  Both matched by hand.
%  The dynamic nearest-marked pattern (Xenia and Tree) is taught here in PROSE;
%  its URL is OWNED by Ch.12 12.3 (/342/E), so it is named, never re-linked.
%  Hand-verified on the running tree T (unweighted, hops): mark {5,13}, query
%  node 6 -> the centroid-ancestor min gives 2 (nearest marked is 5).
% =====================================================================
\section{The Centroid Tree as a Distance Index}\label{sec:c18-index}

\begin{previouslybox}[Ch.~16 answers a distance between two nodes; this is a distance to a \emph{set}]
Chapter~16 (\chref{c11}) turned \(\dist(u,v)\) into an \(\bigO{1}\) lookup once
depths and \(\lca\) are precomputed. But a pairwise oracle answers \emph{one} pair per
call. Two tasks recur and are not pairwise: ``\emph{mark} a node, then ask the distance
to the \emph{nearest} marked node'' (an online set that grows), and ``\emph{count} the
pairs within a distance budget'' (a whole-component sweep). Chapter~22 (\chref{c16})
handles a fixed path's \emph{range}; it does not index a dynamic set and cannot
enumerate all pairs. The centroid tree does: each centroid \(c\) keeps a distance table
for its whole component, and a node \(x\) sits in exactly the \(\bigO{\log n}\) tables
of its centroid ancestors (\(\thmref{c18-height}\)) --- so \emph{touching \(x\) once
touches \(\log n\) component-wide tables}, which is the index this section builds.
\end{previouslybox}

\begin{topicintro}
Root the view not at a vertex of \(T\) but at a centroid \(c\) and record, for every
node \(x\) of \(c\)'s live component, the single number \(d_c(x)=\dist(c,x)\). Do that
at every centroid of the decomposition and you own a family of distance tables with a
special property: for any node \(x\), the tables containing \(x\) are precisely those of
\(x\)'s ancestors in the centroid tree --- at most \(\lceil\log_2 n\rceil\) of them.
That is enough to answer two very different questions. \emph{Online:} to mark \(x\),
write \(d_c(x)\) into each ancestor table's running best; to ask the nearest marked node
to \(v\), take \(\min_c\big(\dist(v,c)+\text{best}[c]\big)\) over \(v\)'s ancestors ---
correct because the highest shared centroid on the path to the true nearest witness
realises equality in the triangle inequality (\(\thmref{c18-nn}\)). \emph{Offline:} to
count pairs at distance within a budget, sweep each centroid's table and use
\(\dist(u,v)=d_c(u)+d_c(v)\) for the pairs its arms straddle, then recurse --- this is
CF~293E, where the budget is \emph{two}-sided (\(\le l\) edges and \(\le w\) weight),
turning the \(\secref{c18-count}\) histogram into a sorted table plus a Fenwick over
weights.

\textbf{Prerequisites.} \(\secref{c18-tree}\) (the decomposition, \(\mathit{par}_c\),
\(\mathit{cd}\), and \(\thmref{c18-height}\)'s \(\bigO{\log n}\) bound) and
\(\secref{c18-count}\) (the top-centroid claim and the same-arm correction). The
nearest-marked query is the pattern the classic \emph{Xenia and Tree} problem needs; its
statement lives in Chapter~12 (\chref{c09}) and is named here, not re-linked.

\textbf{After this section you can}
\begin{itemize}[nosep]
  \item store, for each node, its list of centroid ancestors with hop-distances, in \(\bigO{n\log n}\); \hfill\lv{4}
  \item answer \emph{mark} and \emph{nearest-marked} online in \(\bigO{\log n}\) (each a
        scan of one node's centroid-ancestor list), and prove the min is exact (\(\thmref{c18-nn}\)); \hfill\lv{5}
  \item count pairs under a two-sided distance budget by sorting one metric and a Fenwick
        on the other, inside each centroid bucket; \hfill\lv{5}
  \item solve CF~293E end to end and hand-check it against both official samples. \hfill\lv{5}
\end{itemize}
\end{topicintro}

% ---------------------------------------------------------------
\subsection{One distance table per centroid}
% ---------------------------------------------------------------

\begin{definition}{The centroid distance index}{c18-index}
Run the decomposition of \(\secref{c18-tree}\). For each centroid \(c\), the
\emph{distance table} of \(c\) is the map \(x \mapsto d_c(x)=\dist(c,x)\) defined for
every \(x\) in \(c\)'s live component at the moment \(c\) was chosen. Store it
implicitly: for every node \(x\), keep the list
\(\mathit{up}(x)=\big[(c_0,d_{c_0}(x)),(c_1,d_{c_1}(x)),\dots\big]\) of \(x\)'s centroid
ancestors \(c_0=x\)'s own centroid, up to the top centroid, each paired with its
hop-distance. Because a centroid is deleted before recursion, \(\mathit{up}(x)\) has one
entry per centroid-tree level, so \(|\mathit{up}(x)|\le\lceil\log_2 n\rceil+1\).
\end{definition}

\noindent\textbf{In plain words.} Every node remembers \emph{which} centroids owned it
and \emph{how far} it sat from each. That short list --- logarithmic, fixed at build
time --- is the entire index. Reading it lets one node talk to whole components at once;
it is what turns ``distance from a set'' or ``count within a radius'' from \(O(n)\) per
call into \(O(\log n)\).

\begin{insightbox}[Why the top shared centroid is the one that answers]
Fix a query node \(v\) and any target node \(m\). Let \(z\) be the highest centroid
(smallest \(\mathit{cd}\)) lying on the path \(v\text{--}m\). At the moment \(z\) was
picked, both \(v\) and \(m\) were in its live component, so \(z\in\mathit{up}(v)\cap
\mathit{up}(m)\) and \(z\) is on the \(v\text{--}m\) path, giving
\(\dist(v,m)=\dist(v,z)+\dist(z,m)=d_z(v)+d_z(m)\) \emph{exactly}. For every other
centroid \(c\) that is an ancestor of \(v\), the triangle inequality gives
\(d_c(v)+d_c(m)\ge\dist(v,m)\). So a \(\min\) over \(v\)'s ancestors of
\(d_c(v)+\text{best}[c]\) (where \(\text{best}[c]\le d_c(m)\)) is never \emph{below} the
true nearest distance and is \emph{attained} at \(z\): the min is exact. This is the whole
magic --- the same list that indexes distances also witnesses them.
\end{insightbox}

\begin{theorem}{Nearest-marked via centroid ancestors is exact}{c18-nn}
Maintain \(\text{best}[c]=\min\{d_c(m): m \text{ marked}\}\) (initially \(+\infty\)) for
every centroid \(c\), updating it on a \emph{mark} at \(m\) by scanning \(\mathit{up}(m)\).
Then \(\min_{(c,\,d_c(v))\in\mathit{up}(v)}\big(d_c(v)+\text{best}[c]\big)\) equals
\(\dist(v,m^\ast)\) for the nearest marked node \(m^\ast\).
\end{theorem}
\begin{proofbox}
Every value in the min is \(\ge\dist(v,m^\ast)\) because \(\text{best}[c]=d_c(m)\) for
some marked \(m\), and \(d_c(v)+d_c(m)\ge\dist(v,m)\ge\dist(v,m^\ast)\). Conversely let
\(z\) be the highest centroid on the path \(v\text{--}m^\ast\). Marking \(m^\ast\) set
\(\text{best}[z]\le d_z(m^\ast)\), and \(z\in\mathit{up}(v)\); the insight box shows the
entry at \(z\) is \(d_z(v)+d_z(m^\ast)=\dist(v,m^\ast)\). The min is therefore both
\(\ge\) and \(\le\dist(v,m^\ast)\), hence equal. \(\checkmark\)
\end{proofbox}

\begin{lstlisting}[caption={Build the centroid distance index and answer mark / nearest-marked online. Each operation scans one node's centroid-ancestor list, so it costs O(log n). This is the Xenia-and-Tree pattern; dfs size and find centroid are exactly lst:c18-build. ASCII only}, label={lst:c18-index}]
const int MAXN = 100005, INF = 1e9;
vector<int> adj[MAXN];
vector<pair<int,int>> up[MAXN];    // (centroid ancestor, hop distance node -> it)
bool dead[MAXN]; int sz[MAXN];
int best[MAXN];                    // best[c] = min dist from centroid c to a marked node

int dfs_size(int u, int p){ sz[u]=1;
    for(int v: adj[u]) if(v!=p && !dead[v]) sz[u]+=dfs_size(v,u);
    return sz[u]; }
int find_centroid(int u, int p, int tot){        // Ch.12 walk-to-heavy, lst:c18-build
    bool moved=true;
    while(moved){ moved=false;
        for(int v: adj[u]) if(v!=p && !dead[v] && sz[v]>tot/2){ p=u; u=v; moved=true; break; } }
    return u; }
void collect_anc(int u, int p, int c, int d){    // record centroid c at distance d in u
    up[u].push_back({c, d});
    for(int v: adj[u]) if(v!=p && !dead[v]) collect_anc(v, u, c, d+1);
}
void build(int entry){
    int tot = dfs_size(entry, -1);
    int c   = find_centroid(entry, -1, tot);
    collect_anc(c, -1, c, 0);                    // distances from c to its live component
    dead[c] = true;                              // wall off c before recursing
    for(int v: adj[c]) if(!dead[v]) build(v);
}
// call  build(1);  then set best[c]=INF for all c.
void mark(int v){                                // O(log n): touch v's ancestor tables
    for(auto &a : up[v]) best[a.first] = min(best[a.first], a.second);
}
int query(int v){                                // O(log n): nearest marked distance
    int r = INF;
    for(auto &a : up[v]) r = min(r, best[a.first] + a.second);
    return r;                                    // exact by thm:c18-nn
}
\end{lstlisting}

\begin{tracebox}[The index on the running tree \(T\) -- mark \(\{5,13\}\), then query node \(6\)]
Ancestor lists (centroid, hop-distance) from the decomposition of \(T\):
\(\mathit{up}(5)=\{(4,2),(1,1),(5,0)\}\), \(\mathit{up}(13)=\{(4,3),(3,2),(9,1),(13,0)\}\),
\(\mathit{up}(6)=\{(4,1),(1,1),(6,0)\}\).\\
\textbf{mark 5:} \(\text{best}[4]=2,\ \text{best}[1]=1,\ \text{best}[5]=0\).\\
\textbf{mark 13:} \(3\not<2\), so \(\text{best}[4]\) stays \(2\); set \(\text{best}[3]=2,\ 
\text{best}[9]=1,\ \text{best}[13]=0\).\\
\textbf{query 6} reads \(\mathit{up}(6)\): at \(4\), \(1+\text{best}[4]=1+2=3\); at \(1\),
\(1+\text{best}[1]=1+1=\mathbf{2}\); at \(6\), \(\text{best}[6]=\infty\). The min is
\(\mathbf{2}\). Ground truth: \(\dist(6,5)=2\) via \(6\text{--}1\text{--}5\), and node
\(5\) is marked, so the nearest marked node to \(6\) is \(5\) at distance \(2\). The tight
entry is centroid \(1\), the top centroid on the path \(6\text{--}5\) --- exactly
\(\thmref{c18-nn}\)'s witness. \(\checkmark\)
\end{tracebox}

% ---------------------------------------------------------------
\subsection{Problem Lab: CF 293E -- counting pairs within a two-sided budget}
% ---------------------------------------------------------------

\solvedproblem{c18-293e}{Close Vertices}{CF 293E}
  {https://codeforces.com/problemset/problem/293/E}{CF 2700 L5}
  {centroid decomposition; a per-centroid (edges, weight) table; sort-by-edges with a Fenwick over weights; same-arm subtraction}

\begin{problemstatement}[Close Vertices (statement and samples verified against the live problem)]
A weighted tree on \(n\) vertices (\(1\le n\le 10^5\)) has a non-negative weight on each
edge (\(\le 10^4\)). The \emph{length} of a path is its edge \emph{count}; its
\emph{weight} is the sum of edge weights. Two vertices are \emph{close} if their unique
path has length at most \(l\) (\(1\le l\le n\)) and weight at most \(w\)
(\(0\le w\le 10^9\)). Count unordered pairs \((u,v)\), \(u<v\), that are close. Edges
are given parent-first: line \(i\) gives \(p_i, w_i\), the edge joining vertex \(i+1\) to
\(p_i\) with weight \(w_i\).

\textbf{Examples (official).}
\begin{verbatim}
Input:            Output:
4 4 6             4
1 3
1 4
1 3

Input:            Output:
6 2 17            9
1 3
2 5
2 13
1 6
5 9
\end{verbatim}
\end{problemstatement}

\begin{attemptgate}[Attempt it yourself first -- \lv{5}, give yourself 40 minutes]
Everything is still the \(\secref{c18-count}\) claim: charge each pair to its top
centroid, where the path splits into two arms. \peek{what two numbers per node does a
top-centroid pair combine?} \nudge{for a node \(x\) at the centroid record
\((a_x,b_x)=(\text{edges from }c\text{ to }x,\ \text{weight from }c\text{ to }x)\); a
straddling pair is close iff \(a_x+a_y\le l\) and \(b_x+b_y\le w\).}
\fullhint{Count, over the whole component's table, the pairs meeting BOTH bounds, then
subtract the same count restricted to each single arm (the same-arm correction of
\(\secref{c18-count}\)). The two-sided count of \((a,b)\) pairs with \(a_i+a_j\le l\) and
\(b_i+b_j\le w\) is a 2-D dominance problem: sort by \(a\), and answer with a Fenwick over
compressed \(b\) values. Prune any arm-walk the moment \(a>l\) or \(b>w\). Recurse; total
\(\bigO{n\log^2 n}\).}
\end{attemptgate}

\solstep{Understanding}
Two independent metrics bound the same path: a hop count and a weight sum. Chapter~22's
HLD counts pairs on \emph{one} chosen path; here we range over \emph{all} pairs, and each
pair must clear \emph{both} thresholds. That kills the plain distance histogram of
\(\secref{c18-count}\) --- it only tracked one integer. The fix keeps the centroid
framing (top-centroid claim, same-arm correction) but replaces the \(O(1)\) histogram
lookup with an ordered structure that answers ``how many stored \(b\) values are \(\le
B\)''.

\solstep{Planning}
At centroid \(c\), for every live node \(x\) store \((a_x,b_x)\). A pair is claimed by
\(c\) iff it straddles two different arms (or one endpoint is \(c\)); then
\(\dist=c\)-sums apply: \(a_x+a_y\le l\), \(b_x+b_y\le w\). By inclusion--exclusion,
\(\text{close pairs at }c=\mathit{count}(\text{all})-\sum_{\text{arm}}\mathit{count}(\text{arm})\),
where \(\mathit{count}(S)\) counts unordered pairs from the multiset \(S\) clearing both
bounds. Then mark \(c\) dead and recurse.

\solstep{Implementing}
\(\mathit{count}(S)\): sort \(S\) by \(a\). Each element \(i\) asks ``how many \(j\in S\)
have \(a_j\le l-a_i\) and \(b_j\le w-b_i\)'' --- a prefix in \(a\) combined with a prefix
in \(b\), i.e.\ 2-D dominance. Process these queries in increasing \(l-a_i\), inserting
each point's \(b\) into a Fenwick as its \(a\) becomes eligible; a query reads a Fenwick
prefix sum. This counts \emph{ordered} pairs including self-pairs; subtract the self-count
and halve. Compress the \(b\)-values of \(S\) (weights up to \(10^9\)).

\begin{lstlisting}[caption={CF 293E: count close pairs. At each centroid, build a per-node (edges, weight) table, count pairs clearing both budgets with a sort plus a Fenwick over compressed weights, and subtract the same-arm counts. dfs size and find centroid match lst:c18-build. ASCII only}, label={lst:c18-293e}]
#include <bits/stdc++.h>
using namespace std;
typedef long long ll;
const int MAXN = 100005;
int n; ll L, W;                       // L = max edges, W = max weight
vector<pair<int,int>> adj[MAXN];      // (neighbour, weight)
bool dead[MAXN]; int sz[MAXN];
ll ans = 0;

// unordered pairs i<j in pts with a_i+a_j<=L and b_i+b_j<=W
ll countPairs(vector<pair<ll,ll>>& pts){
    int m = pts.size(); if(m<=1) return 0;
    vector<ll> bs; for(auto&t:pts) bs.push_back(t.second);
    sort(bs.begin(),bs.end()); bs.erase(unique(bs.begin(),bs.end()),bs.end());
    int B = bs.size(); vector<int> bit(B+1,0);
    auto add=[&](int i,int v){ for(;i<=B;i+=i&-i) bit[i]+=v; };
    auto qry=[&](int i){ int s=0; for(;i>0;i-=i&-i)s+=bit[i]; return s; };
    sort(pts.begin(),pts.end());                       // by a ascending
    vector<pair<ll,int>> q(m);                         // (L - a_i, i)
    for(int i=0;i<m;i++) q[i]={L-pts[i].first, i};
    sort(q.begin(),q.end());
    ll ordered=0; int p=0;
    for(auto&qu:q){
        while(p<m && pts[p].first<=qu.first){          // a_j <= L - a_i
            int bidx=lower_bound(bs.begin(),bs.end(),pts[p].second)-bs.begin()+1;
            add(bidx,1); p++;
        }
        int bi=upper_bound(bs.begin(),bs.end(),W-pts[qu.second].second)-bs.begin();
        ordered += qry(bi);                            // b_j <= W - b_i
    }
    ll self=0;                                         // remove (i,i) matches
    for(auto&t:pts) if(2*t.first<=L && 2*t.second<=W) self++;
    return (ordered-self)/2;                           // ordered -> unordered
}
void collect(int u,int p,ll a,ll b,vector<pair<ll,ll>>& arm,vector<pair<ll,ll>>& all){
    if(a>L || b>W) return;                             // cannot clear a bound through c
    arm.push_back({a,b}); all.push_back({a,b});
    for(auto&e:adj[u]) if(e.first!=p && !dead[e.first])
        collect(e.first,u,a+1,b+e.second,arm,all);
}
ll process(int c){
    vector<pair<ll,ll>> all; all.push_back({0,0});     // the centroid itself
    vector<vector<pair<ll,ll>>> arms;
    for(auto&e:adj[c]){ if(dead[e.first]) continue;
        vector<pair<ll,ll>> arm;
        collect(e.first,c,1,e.second,arm,all);
        arms.push_back(arm);
    }
    ll tot = countPairs(all);
    for(auto&arm:arms) tot -= countPairs(arm);         // same-arm correction
    return tot;
}
int dfs_size(int u,int p){ sz[u]=1;
    for(auto&e:adj[u]) if(e.first!=p && !dead[e.first]) sz[u]+=dfs_size(e.first,u);
    return sz[u]; }
int find_centroid(int u,int p,int tot){
    bool moved=true;
    while(moved){ moved=false;
        for(auto&e:adj[u]) if(e.first!=p && !dead[e.first] && sz[e.first]>tot/2){ p=u;u=e.first;moved=true;break; } }
    return u; }
void decompose(int entry){
    int tot=dfs_size(entry,-1); int c=find_centroid(entry,-1,tot);
    ans += process(c);
    dead[c]=true;
    for(auto&e:adj[c]) if(!dead[e.first]) decompose(e.first);
}
int main(){
    scanf("%d %lld %lld",&n,&L,&W);
    for(int i=2;i<=n;i++){ int p,w; scanf("%d %d",&p,&w);
        adj[i].push_back({p,w}); adj[p].push_back({i,w}); }
    decompose(1);
    printf("%lld\n",ans);
}
\end{lstlisting}

\begin{tracebox}[Sample 2 by hand -- \(n=6,\ l=2,\ w=17\), tree \(1\!-\!2(3),2\!-\!3(5),2\!-\!4(13),1\!-\!5(6),5\!-\!6(9)\)]
\textbf{Top centroid \(1\)} (removing \(1\) gives pieces \(\{2,3,4\}\) size 3 and
\(\{5,6\}\) size 2, max \(3\le 3\)). Table \((a,b)\) at \(1\): \(1(0,0)\), \(2(1,3)\),
\(3(2,8)\), \(4(2,16)\), \(5(1,6)\), \(6(2,15)\). Count \(\mathit{all}\) with \(a_i+a_j\le2\),
\(b_i+b_j\le17\): with \(1\): \((1,2),(1,3),(1,4),(1,5),(1,6)\) all clear (5); across arms
\(2\text{--}5\): \(a\,1{+}1{=}2,\ b\,3{+}6{=}9\) clear (1); \(2\text{--}6\) and \(3,4\)
with \(5,6\) fail the edge bound (\(a>2\)); \(3\text{--}4\): \(b\,8{+}16{=}24\) fails.
So \(\mathit{count}(\text{all})=6\). Arm \(\{2,3,4\}\): every pair has \(a\ge1{+}2=3>2\)
\(\Rightarrow0\). Arm \(\{5,6\}\): \(a\,1{+}2=3>2\Rightarrow0\). Contribution \(6-0-0=6\).\\
\textbf{Recurse arm \(\{2,3,4\}\),} centroid \(2\): table \(2(0,0),3(1,5),4(1,13)\);
\(\mathit{count}(\text{all})\): \((2,3)\) ok, \((2,4)\) ok, \((3,4)\) \(b\,18>17\) fails
\(\Rightarrow2\); arms singleton \(\Rightarrow0\). Contribution \(2\).\\
\textbf{Recurse arm \(\{5,6\}\),} centroid \(5\): \(5(0,0),6(1,9)\); pair clears
\(\Rightarrow1\). Contribution \(1\).\\
\textbf{Total} \(6+2+1=9\). \(\checkmark\) matches the official output.
\end{tracebox}

\begin{complexitybox}[Two-sided distance counting]
Pruned by \(a\le l\) and \(b\le w\), each centroid's table is at most its component size;
\(\mathit{count}(\cdot)\) is \(\bigO{s\log s}\) (a sort and a Fenwick sweep) and runs on
\(\mathit{all}\) plus every arm, so per centroid \(\bigO{s\log s}\). A node is in
\(\bigO{\log n}\) components (\(\thmref{c18-height}\)), giving \(\bigO{n\log^2 n}\) time
overall and \(\bigO{n}\) extra space beyond the tables. The single \(\log n\) of
\(\secref{c18-count}\) became \(\log^2 n\) because the second metric needs an ordered
structure rather than an array index --- the price of a two-sided constraint.
\end{complexitybox}

\begin{warningbox}[Three ways this goes wrong]
\begin{itemize}[nosep]
  \item \textbf{Forgetting the same-arm subtraction.} \(\mathit{count}(\text{all})\) also
        counts pairs that live in one child arm, whose real top centroid is deeper. Those
        pairs are counted here \emph{and} down there --- double count. The
        \(\sum_{\text{arm}}\mathit{count}(\text{arm})\) line is not decoration.
  \item \textbf{Self-pairs.} \(\mathit{count}(\text{all})\) treats \(u=v\): \((c,c)\) and
        any node with \(2a\le l,\ 2b\le w\) sneak in. The \(\mathit{self}\) correction
        removes them before halving; skip it and every answer is too large.
  \item \textbf{Ordering vs unordered.} The dominance sweep counts ordered pairs
        \((i,j)\). Halving is only valid after self-pairs are gone; halve first and you
        floor away real pairs.
\end{itemize}
\end{warningbox}

\begin{hardeningbox}[Where this skeleton generalises]
\begin{itemize}[nosep]
  \item \textbf{One budget, count all.} Drop \(b\) and use an array histogram: this is
        exactly \(\secref{c18-count}\) (\(\coderef{c18-count}\)) --- the Fenwick is the
        histogram upgraded to a second, unordered metric.
  \item \textbf{Farthest instead of nearest.} Replace \(\min\) in \(\coderef{c18-index}\)'s
        query with \(\max\); the same ancestor list reports the farthest marked node (and
        the tree's diameter via \(\chref{c09}\)'s farthest baseline, cited not recomputed).
  \item \textbf{Sum over all, not count.} Change \(\mathit{count}(\cdot)\) to a Fenwick
        that accumulates \(b\)-sums as well as counts and you get ``sum of path weights
        over close pairs'' for free --- the Summit (\(\probref{c18-summit}\)) leans on this.
\end{itemize}
\end{hardeningbox}

\begin{recallbox}[24.3 in one breath]
\emph{Index the component, not the pair.} Each centroid owns a distance table; a node
lives in \(\bigO{\log n}\) of them (\(\defref{c18-index}\)). Mark/write and
nearest/query read that short list, exact by \(\thmref{c18-nn}\)'s top-shared-centroid
witness. Offline counting keeps the top-centroid claim but swaps the histogram for a
sort-plus-Fenwick so a second metric fits --- CF~293E counts pairs within \(\le l\) edges
\emph{and} \(\le w\) weight in \(\bigO{n\log^2 n}\) (\(\coderef{c18-293e}\)).
\end{recallbox}
